\scanpage{189}
% archive sheet number: 188

\unsure{garantit} que $A$, en tant que co-algèbre, soit limite inductive de ``ss-Coalgèbres'', de type fini en tant que modules sur $k$ \deleted{\unsure{s'il en est ainsi}} \deleted{et, de préférence} [et de préférence facteurs directs en tant que modules sur $k$]. \uline{S'il} en est cependant ainsi, alors $\check A$ peut être regardé comme l'algèbre affine d'un ind-(schéma affine) sur $k$, \inserted{[pas nécessairement un schéma formel]} \deleted{\unsure{, qui sera un groupe}} et \deleted{sera}\inserted{définit} \unsure{=} un \uline{groupe} (ind-affine) sur $k$, désigné alors par $D(G)$ et appelé \uline{groupe dual} (au sens Cartier) de $G$. \deleted{\unsure{Ce sera le}}\inserted{Le foncteur} \inserted{de $k$-algèbres} représenté par ce \deleted{\unsure{progr.}} groupe ind.affine sera précisément $\underline{\mathrm{Hom}}_{k\text{-grps}}(G,\mathbb G_m)$.
% slanted
\margin{Si, \unsure{dans le} cas d'un corps cela \unsure{va de bien|va de soi}\unsure{...}}

Si \deleted{$G$} $k$ est \deleted{de type fini} un corps et $G$ de type fini sur $k$, alors si $k$ est de car $p>0$, d'après Cartier $D(G)$ est bien un groupe (ind-fini).
\[
\left\{\begin{aligned}
D(\mathbb G_m)&=\deleted{\unsure{\mathbb Z}}\mathbb Z\\
D(\mathbb Z/\ell\mathbb Z)&\simeq\mathbb Z/\ell\mathbb Z
\quad\text{\deleted{\unsure{naturellement}} si $\ell$ premier \unsure{à car.}, et}\\[-2pt]
&\hspace{28mm}\text{si les racines primitives $\ell$-ièmes de l'unité \unsure{existent}}\\[-2pt]
&\hspace{28mm}\text{dans \deleted{\unsure{$\Omega$}}$k$}\\
D(\mathbb Z/p^n\mathbb Z)&\simeq\mu_{p^n}\quad\text{\deleted{\unsure{par ce}}}\\
D(\mu_{p^n})&\simeq\mathbb Z/p^n\mathbb Z\\
D(\mathbb G_a)&=\Bigl(\text{Algèbre enveloppante \inserted{\uline{\unsure{avec distinction de puiss $p^i$|sans distinction de puiss $p^i$}}}}\\[-2pt]
&\hspace{8mm}\text{de la $p$-algèbre de Lie libre engendrée par un élément}\\[-2pt]
&\hspace{8mm}\text{(sur un corps de car $p>0$)}\Bigr)
\end{aligned}\right.
\]
